\documentclass[10pt]{beamer}
% \documentclass[handout,10pt]{beamer}
\usepackage{my_style}
%%%%%%%%%%%%%%%%%%%%%%%%%%%%%%%%%%%%%%%%%%%%%%%%%%%%%%%%%%%%%%%%%%%%%%%%

\title[Math/CS 375]{Lecture 07 - A}
\subtitle{\S 8.1 LU Factorization}

\begin{document}
% -------------------------------------------------
\begin{frame}
  \titlepage
\end{frame}
% -------------------------------------------------
\begin{frame}{Goals}
We will answer these questions:\\[3mm]
\begin{itemize}
    \item What are efficient strategies for solving $AX=B$, a multiple right-hand side problem?
    \item What is the benefit of factorizing $A=LU$ in terms of solving $Ax=b$?
    \item How is Gaussian elimination expressed in terms of matrix multiplication?
    \item How do you construct an LU factorization.
\end{itemize}
\end{frame}

%%%%%%%%%%%%%%%%%%%%%%%%%%%%%%%%%%%%%%%%%%%%%%%%%%%%%%%%%%%%%%%%%%%%%%%
\sectionframe{Multiple Right Hand Side}
%%%%%%%%%%%%%%%%%%%%%%%%%%%%%%%%%%%%%%%%%%%%%%%%%%%%%%%%%%%%%%%%%%%%%%%%


\begin{frame}{Motivation}

\begin{block}{Single Solve}
Suppose $A\in\mathbb{R}^{n\times n}$. The singular problem
$$
    Ax=b
$$
costs $\mO(n^3)$ to solve.
\end{block}

Many applications require solving $Ax=b_i$ repeatedly
with the \textbf{same matrix $A$}:
\[
    Ax^{(1)}=b^{(1)},\quad
    Ax^{(2)}=b^{(2)},\quad \ldots,\quad
    Ax^{(m)}=b^{(m)}.
\]


\end{frame}

%%%%%%%%%%%%%%%%%%%%%%%%%%%%%%%%%%%%%%%%%%%%%%%%%%%%%%%%%%%%%%%%%%%%%%%%

\begin{frame}{Consider this problem...}

\begin{block}{Equivalent Problems}
\begin{itemize}
    \item Solving $Ax=b$ for \textbf{multiple right hand sides} (RHS), i.e.
    $$ Ax_1 = b_1,\; Ax_2=b_2,\; \dots,\; Ax_m = b_m$$ 
    
    \item Solving for $X$ in the equation 
    $$ \underbrace{A}_{n \times n} \underbrace{X}_{n\times m} = \underbrace{B}_{n\times m}.$$
\end{itemize}
\end{block}

\question{
    What are the strategies to solve $AX=B?$ What will they cost?
}

\end{frame}

%%%%%%%%%%%%%%%%%%%%%%%%%%%%%%%%%%%%%%%%%%%%%%%%%%%%%%%%%%%%%%%%%%%%%%%%

\begin{frame}{Multiple Right Hand Sides}{Options}


\begin{block}{Option A: Do Gaussian elimination (GE) $m$ times }

        \begin{itemize}
            \item Pro: Never compute $A^{-1}$
            \item Con: But it's expensive $\mO(m n^3)$
        \end{itemize}

\end{block}

\onslide<2->{
\begin{block}{Option B: Compute $A^{-1}$ once and reuse it }
\begin{itemize}
            \item Pro: Get to reuse $A^{-1}$ for each RHS
            \item Con: but you're computing $A^{-1}$ 
\end{itemize}
\end{block}
}


\vspace{2mm}

\onslide<3->{
\keyidea{
    \begin{itemize}
        \item Options A and B are \textbf{both bad}
        \item \textbf{Matrix decomposition} combines the best parts of Options A and B
    \end{itemize}
}
}

\end{frame}

\begin{frame}{What is Matrix Decomposition?}

{
\define{Matrix Decomposition}{\textbf{Matrix decomposition} or \textbf{matrix factorization} is a factorization of a matrix $A$ into a product of matrices, e.g. 
$$ A = A_1 A_2 $$}
}

\begin{block}{Three Types}
We will learn three types of decompositions
\begin{itemize}
    \item $A = LU$
    \item $A = LDL^T$
    \item $A = L L^T$
\end{itemize}
\end{block}


\end{frame}

%%%%%%%%%%%%%%%%%%%%%%%%%%%%%%%%%%%%%%%%%%%%%%%%%%%%%%%%%%%%%%%%%%%%%%%%
\sectionframe{LU Solve}
%%%%%%%%%%%%%%%%%%%%%%%%%%%%%%%%%%%%%%%%%%%%%%%%%%%%%%%%%%%%%%%%%%%%%%%%
\begin{frame}{$LU$ Factorization}
\define{LU Factorization}{\textbf{LU factorization} or \textbf{LU decomposition} is a factorization for a $n\times n$ matrix $A$ into a product
\begin{equation*}
    A = LU
\end{equation*}
where $L$ is lower triangular with ones on the diagonal, and $U$ is upper triangular.
}
\begin{equation*}
    L = \begin{bmatrix}   1       &     0      & \cdots &             & 0 \\
                       \ell_{2,1} &     1      &   0    &             & 0 \\
                       \ell_{3,1} & \ell_{3,2} &   1    &             & 0 \\
                          \vdots  &  \vdots    &        & \ddots      & \vdots  \\
                       \ell_{n,1} & \ell_{n,2} & \cdots & \ell_{n,n-1} & 1
        \end{bmatrix} \quad 
    U = \begin{bmatrix} u_{1,1} & u_{1,2} & u_{1,3} & \cdots & u_{1,n} \\
                          0     & u_{2,2} & u_{2,3} & \cdots & u_{2,n} \\
                          0     &   0     & \ddots  & \ddots & \vdots \\
                         \vdots & \vdots  &         &        & u_{n-1,n} \\
                          0     &   0     &         &        & u_{n,n}
        \end{bmatrix}
\end{equation*}
\end{frame}
%%%%%%%%%%%%%%%%%%%%%%%%%%%%%%%%%%%%%%%%%%%%%%%%%%%%%%%%%%%%%%%%%%%%%%%%
\begin{frame}{Main Advantage}

\begin{itemize}
    \item Usually, solving $Ax=b$ costs $\mO(n^3)$, but...
\begin{theorem}
    If $A$ is factorized as $A=LU$ then $Ax=b$ can be solved in $\mO(n^2)$ operations.
\end{theorem}
\vspace{3mm}
        \item If $A=LU$, then it is an \textbf{entire order cheaper} to solve the system
    \end{itemize}
\vspace{3mm}
\question{How does this work?}
\end{frame}
%%%%%%%%%%%%%%%%%%%%%%%%%%%%%%%%%%%%%%%%%%%%%%%%%%%%%%%%%%%%%%%%%%%%%%%%
\begin{frame}{The $LU$ Solve}
\textbf{Steps to solve $Ax=b$:}\\[4mm]

\begin{enumerate}
    \item Substitute $A=LU$ in $Ax=b$ to get 
        $$ LUx=b$$

    \item Set $y=Ux$ so that 
        $$ Ly = b$$
    \item Use forward substitution $(\mO(n^2))$ to compute 
        $$ y = L^{-1}b$$
    \item Use backward substitution $(\mO(n^2))$ to compute 
        $$ x = U^{-1}y$$
\end{enumerate}
\end{frame}
%%%%%%%%%%%%%%%%%%%%%%%%%%%%%%%%%%%%%%%%%%%%%%%%%%%%%%%%%%%%%%%%%%%%%%%%
\begin{frame}{Revisit Options for Solving $AX=Bs$ }


\begin{block}{Option A: Do Gaussian elimination (GE) $m$ times }

        \begin{itemize}
            \item Pro: Never compute $A^{-1}$
            \item Con: But it's expensive $\mO(m n^3)$
        \end{itemize}

\end{block}


\begin{block}{Option B: Compute $A^{-1}$ once and reuse it }
\begin{itemize}
            \item Pro: Get to reuse $A^{-1}$ for each RHS
            \item Con: but you're computing $A^{-1}$ 
\end{itemize}
\end{block}


\begin{block}{Option C: Factorize $A=LU$ once and reuse it}
        \begin{itemize}
            \item $A^{-1}$ is never computed
            \item Cheaper: $\mO(n^3 + mn^2)$
        \end{itemize}    

\end{block}

    \keyidea{
        LU factorization is safer and faster than both options 
    }
\end{frame}

%%%%%%%%%%%%%%%%%%%%%%%%%%%%%%%%%%%%%%%%%%%%%%%%%%%%%%%%%%%%%%%%%%%%%%%%
\begin{frame}{$LU$ Solve Method}

\method{$LU$ Solve}
{Factorization $A=LU$, RHS $b$}
{Solution to $Ax=b$}
{The LU solve costs only $\mO(n^2)$ which is cheaper than GE.}
\vspace{3mm}
Comments:
\begin{itemize}
    \item The triangular solves are both cost $\mO(n^2)$
    \item The LU solve costs around $2n^2$ FLOPs
    \item Especially useful when solving multiple RHS, $m\gg1$
\end{itemize}


\onslide<2>{
\question{So how do you find $A=LU$?}
}
\end{frame}


%%%%%%%%%%%%%%%%%%%%%%%%%%%%%%%%%%%%%%%%%%%%%%%%%%%%%%%%%%%%%%%%%%%%%%%%
\sectionframe{Elimination Matrices}
%%%%%%%%%%%%%%%%%%%%%%%%%%%%%%%%%%%%%%%%%%%%%%%%%%%%%%%%%%%%%%%%%%%%%%%

\begin{frame}[shrink]{Consider This}
    \begin{example}
        Calculate the following matrix vector product
            
            $$  Mv =     
            \begin{bmatrix}
          1 & 0 & 0\\
        -2 & 1 & 0\\
        1 & 0 & 1\\
        \end{bmatrix} \begin{bmatrix}2 \\ 4 \\ -2\end{bmatrix}$$ 
    \end{example}

\begin{itemize}
    \item Product is 
        \[
        \begin{bmatrix}
          1 & 0 & 0\\
        -2 & 1 & 0\\
        1 & 0 & 1\\
        \end{bmatrix}
        \begin{bmatrix}
        2 \\ 4 \\ -2
        \end{bmatrix}  =
        \begin{bmatrix}
        2 \\ 0 \\ 0
        \end{bmatrix}
        \]
    \item Think of $M$ as an \emph{operator} that zeros out certain elements in $v$
    $$
    \begin{bmatrix}2 \\ 4 \\ -2\end{bmatrix} \longrightarrow          \begin{bmatrix}
        2 \\ 0 \\ 0
        \end{bmatrix}
    $$
    \item $M$ is an example of an \textbf{elimination matrix}
\end{itemize}
\end{frame}

%%%%%%%%%%%%%%%%%%%%%%%%%%%%

\begin{frame}{Understanding the Operator}
    
     $$\begin{bmatrix}
          1 & 0 & 0\\
        -2 & 1 & 0\\
        1 & 0 & 1\\
        \end{bmatrix} \begin{bmatrix}2 \\ 4 \\ -2\end{bmatrix} =
        \begin{bmatrix}
        2 \\ 0 \\ 0
        \end{bmatrix}$$ 

    \begin{block}{Top Row}
         Top row $\begin{bmatrix}
                1 & 0 & 0
            \end{bmatrix}$ says leave the 2 alone.\\
            The 2 remains a 2.
    \end{block}
    
    \begin{block}{Second Row}
            Second row $\begin{bmatrix}
                -2 & 1 & 0
            \end{bmatrix}$ says to add -2 copies of 2 to the 4\\
            The 4 becomes $4-2*2 = 0$.
    \end{block}

       \begin{block}{Third Row}
            Third row $\begin{bmatrix}
                1 & 0 & 1
            \end{bmatrix}$ says to add 1 copy of 2 to the -2.\\
            The -2 becomes $-2+1*2 = 0$.
    \end{block}
\end{frame}
%%%%%%%%%%%%%%%%%%%%%%%%%%%%

\begin{frame}{Understanding the Operator}
    
     $$\begin{bmatrix}
          1 & 0 & 0\\
        -2 & 1 & 0\\
        1 & 0 & 1\\
        \end{bmatrix} \begin{bmatrix}2 \\ 4 \\ -2\end{bmatrix} =
        \begin{bmatrix}
        2 \\ 0 \\ 0
        \end{bmatrix}$$ 

    \begin{block}{Top Row}
         Top row $\begin{bmatrix}
                1 & 0 & 0
            \end{bmatrix}$ says leave the 2 alone.\\
            The 2 remains a 2.
    \end{block}
    
    \begin{block}{Second Row}
            Second row $\begin{bmatrix}
                -2 & 1 & 0
            \end{bmatrix}$ says to add -2 copies of 2 to the 4\\
            The 4 becomes $4-2*2 = 0$.
    \end{block}

       \begin{block}{Third Row}
            Third row $\begin{bmatrix}
                1 & 0 & 1
            \end{bmatrix}$ says to add 1 copy of 2 to the -2.\\
            The -2 becomes $-2+1*2 = 0$.
    \end{block}
\end{frame}

\begin{frame}{What about a Matrix?}

    \begin{itemize}
        \item This matrix is able to eliminate elements under the top element of $v$
     $$Mv=\begin{bmatrix}
          1 & 0 & 0\\
        -2 & 1 & 0\\
        1 & 0 & 1\\
        \end{bmatrix} \begin{bmatrix}2 \\ 4 \\ -2\end{bmatrix} =
        \begin{bmatrix}
        2 \\ 0 \\ 0
        \end{bmatrix}$$ 
        \onslide<2>{
        \question{
        What is the effect of multiplying $M$ to another matrix?
            $$  MA =     
            \begin{bmatrix}
          1 & 0 & 0\\
        -2 & 1 & 0\\
        1 & 0 & 1\\
        \end{bmatrix} \begin{bmatrix}2 &1&1\\ 4 &3&2\\ -2 &-1&5\end{bmatrix}$$ 
        }

        }
     \end{itemize}
\end{frame}


%%%%%%%%%%%%%%%%%%%%%%%%%%%%

\begin{frame}[shrink]{What about a Matrix?}
    \begin{example}
        Calculate the following matrix  product
            
            $$  MA =     
            \begin{bmatrix}
          1 & 0 & 0\\
        -2 & 1 & 0\\
        1 & 0 & 1\\
        \end{bmatrix} \begin{bmatrix}2 &1&1\\ 4 &3&2\\ -2 &-1&5\end{bmatrix}$$ 
    \end{example}

\begin{itemize}
    \item First row: do nothing
        \[
        \begin{bmatrix}
          1 & 0 & 0\\
        -2 & 1 & 0\\
        1 & 0 & 1\\
        \end{bmatrix}
         \begin{bmatrix}2 &1&1\\ 4 &3&2\\ -2 &-1&5\end{bmatrix}  =
        \begin{bmatrix}
        2 & 1 & 1\\ \phantom0 &\phantom1 & \phantom1\\ \phantom0&\phantom0&\phantom6
        \end{bmatrix}
        \]
          \item<2-> Second row: add -2 copies of $R_1$ to $R_2$
        \[
        \begin{bmatrix}
          1 & 0 & 0\\
        -2 & 1 & 0\\
        1 & 0 & 1\\
        \end{bmatrix}
         \begin{bmatrix}2 &1&1\\ 4 &3&2\\ -2 &-1&5\end{bmatrix}  =
        \begin{bmatrix}
        2 &1 & 1\\ 0 &1 & 0\\ \phantom0&\phantom0&\phantom6
        \end{bmatrix}
        \]
             \item<3-> Third row: add 1 copy of $R_1$ to $R_3$
        \[
        \begin{bmatrix}
          1 & 0 & 0\\
        -2 & 1 & 0\\
        1 & 0 & 1\\
        \end{bmatrix}
         \begin{bmatrix}2 &1&1\\ 4 &3&2\\ -2 &-1&5\end{bmatrix}  =
        \begin{bmatrix}
        2 &1 & 1\\ 0 &1 & 0\\ 0&0&6
        \end{bmatrix}
        \]
\end{itemize}
\end{frame}

%%%%%%%%%%%%%%%%%%%%%%%%%%%%

\begin{frame}{Elimination Matrix $\times$ Matrix}
        \question{
        What happens when you multiply an elimination matrix to another matrix?
            $$  MA =     
            \begin{bmatrix}
          1 & 0 & 0\\
        -2 & 1 & 0\\
        1 & 0 & 1\\
        \end{bmatrix} \begin{bmatrix}2 &1&1\\ 4 &3&2\\ -2 &-1&5\end{bmatrix}$$ 
        } 

    \begin{itemize}
        \item<1-> The first row is preserved 
        \[
        \begin{bmatrix}
          1 & 0 & 0\\
        -2 & 1 & 0\\
        1 & 0 & 1\\
        \end{bmatrix}
         \begin{bmatrix}2 &1&1\\ 4 &3&2\\ -2 &-1&5\end{bmatrix}  =
        \begin{bmatrix}
        2 & 1 & 1\\ \phantom0 &\phantom1 & \phantom1\\ \phantom0&\phantom0&\phantom6
        \end{bmatrix}
        \]
        \item<2-> Elements under $a_{11}$ are eliminated 
        \[
        \begin{bmatrix}
          1 & 0 & 0\\
        -2 & 1 & 0\\
        1 & 0 & 1\\
        \end{bmatrix}
         \begin{bmatrix}2 &1&1\\ 4 &3&2\\ -2 &-1&5\end{bmatrix}  =
        \begin{bmatrix}
        \boxed{2} &1 & 1\\ \red{0} &1 & 0\\ \red{0} &0&6
        \end{bmatrix}
        \]
    \end{itemize}
\end{frame}

%%%%%%%%%%%%%%%%%%%%%%%%%%%%%%%%%%%%%%%%%%%%%%%%%%%%%%%%%%%%%%%%%%%%%%%%
\begin{frame}{Designing the Elimination Matrix}
\begin{example}
    
    Find a matrix $4\times 4$ elimination matrix $M$ that eliminates the elements under $a_{11}=6$ when multiplied as $MA$, where
    $A = \begin{bmatrix}
    6 & -2 & 2 & 4\\
    12 & -8 & 6 & 10\\
    3 & -13 & 9 & 3\\
    -6 & 4 & 1 & -18\\
    \end{bmatrix}$
\end{example}
\begin{itemize}

    \item Formula:  
    $\begin{bmatrix}
    1 & 0 & 0 & 0\\
    -m_2 & 1 & 0 & 0\\
    -m_3 & 0 & 1 & 0\\
    -m_4 & 0 & 0 & 1\\
    \end{bmatrix}$, where $m_2= \frac{a_{21}}{a_{11}}$, $m_3=\frac{a_{31}}{a_{11}}$ and $m_4=\frac{a_{41}}{a_{11}}$
    \item Thus our elimination matrix is 
    $$
    \begin{bmatrix}
    1 & 0 & 0 & 0\\
    -2 & 1 & 0 & 0\\
    -\frac{1}{2} & 0 & 1 & 0\\
    1 & 0 & 0 & 1\\
    \end{bmatrix}
    \begin{bmatrix}
    \boxed{6} & -2 & 2 & 4\\
    12 & -8 & 6 & 10\\
    3 & -13 & 9 & 3\\
    -6 & 4 & 1 & -18\\
    \end{bmatrix} = 
    \begin{bmatrix}
    6 & -2 & 2 & 4\\
    0 & -4 & 2 & 2\\
    0 & -12 & 8 & 1\\
    0 & 2 & 3 & -14\\
    \end{bmatrix}
    $$
\end{itemize}
\end{frame}

% \begin{frame}[shrink]{Recall: Gaussian Elimination Step}
% \begin{itemize}
%     \item Eliminate elements under the pivot element in the first column.
%     \item $x'$ indicates a value that has been changed once.
% \begin{center}
% 	\renewcommand{\arraystretch}{1.3}
% 	\footnotesize
% 	\begin{tabular}{ccccc}
% 		$\begin{bmatrix}\boxed{x} & x & x & x & x \\
% 						   x  & x  & x  & x  & x  \\
% 						   x  & x  & x  & x  & x \\
% 						   x  & x  & x  & x  & x \end{bmatrix}$
% 		& $\longrightarrow$ &
% 		$\begin{bmatrix}\boxed{x} & x & x & x & x \\
% 						   0  & x' & x' & x' & x' \\
% 						   x  & x  & x  & x  & x \\
% 						   x  & x  & x  & x  & x \end{bmatrix}$ \\[20pt]
% 		& $\longrightarrow$ &
% 		$\begin{bmatrix}\boxed{x} & x & x & x & x \\
% 						   0  & x' & x' & x' & x' \\
% 						   0  & x' & x' & x' & x' \\
% 						   x  & x  & x  & x  & x \end{bmatrix}$ \\[20pt]
% 		& $\longrightarrow$ &
% 		$\begin{bmatrix}\boxed{x} & x & x & x & x \\
% 						   0  & x' & x' & x' & x' \\
% 						   0  & x' & x' & x' & x' \\
% 						   0  & x' & x' & x' & x'\end{bmatrix}$ &
% 	\end{tabular}
% \end{center}
%     \item You need to calculate the \textbf{multiplier} $m_k = a_{k1}/a_{11}$, $k>1$
% \end{itemize}
% \end{frame}

%%%%%%%%%%%%%%%%%%%%%%%%%%%%%%%%%%%%%%%%%%%%%%%%%%%%%%%%%%%%%%%%%%%%%%%%
\sectionframe{Constructing $U$}
%%%%%%%%%%%%%%%%%%%%%%%%%%%%%%%%%%%%%%%%%%%%%%%%%%%%%%%%%%%%%%%%%%%%%%%%

\begin{frame}{Goal}

\begin{itemize}
    \item We want to construct the $U$ in $$A=LU$$ 

    \item Multiply both sides by $L^{-1}$
        \begin{align*} 
        L^{-1}A&=L^{-1}LU\\
        L^{-1}A&=U
        \end{align*}

    \begin{theorem}
        If $L$ is lower triangular then $L^{-1}$ is also lower triangular.
    \end{theorem}
\end{itemize}

\keyidea{
This equation
$$
L^{-1}A=U
$$
tells us that we construct $U$ by multiplying $A$ by some \textit{lower triangular matrix.}
}
\end{frame}

%%%%%%%%%%%%%

\begin{frame}{What is the lower triangular matrix?}
    \begin{itemize}
        \item We know that elimination matrices are lower triangular 
        $$
        M_1=    \begin{bmatrix}
    1 & 0 & 0 & 0\\
    -2 & 1 & 0 & 0\\
    -\frac{1}{2} & 0 & 1 & 0\\
    1 & 0 & 0 & 1\\
    \end{bmatrix}
    $$

    \item<2-> But the multiplication doesn't always create an upper triangular 
    $$\underbrace{
    \begin{bmatrix}
    1 & 0 & 0 & 0\\
    -2 & 1 & 0 & 0\\
    -\frac{1}{2} & 0 & 1 & 0\\
    1 & 0 & 0 & 1\\
    \end{bmatrix}}_{M_1}
    \underbrace{\begin{bmatrix}
    \boxed{6} & -2 & 2 & 4\\
    12 & -8 & 6 & 10\\
    3 & -13 & 9 & 3\\
    -6 & 4 & 1 & -18\\
    \end{bmatrix}}_{A} = 
    \underbrace{
    \begin{bmatrix}
    6 & -2 & 2 & 4\\
    0 & -4 & 2 & 2\\
    0 & -12 & 8 & 1\\
    0 & 2 & 3 & -14\\
    \end{bmatrix}}_{\text{not upper triangular :(}}
    $$
    \item<3-> What can we do?
    \end{itemize}
\end{frame}

%%%%%%%%%%

\begin{frame}[shrink]{Use Multiple Elimination Matrices}

\begin{center}
	\renewcommand{\arraystretch}{1.3}
	\footnotesize
	\begin{tabular}{ccccc}
		$\underbrace{\begin{bmatrix}\boxed{x} & x & x & x   \\
						   x  & x  & x  & x   \\
						   x  & x  & x  & x  \\
						   x  & x  & x  & x  \end{bmatrix}}_{A}$
		& $\longrightarrow$ &
		$\begin{bmatrix}\boxed{x} & x & x & x    \\
						   0  & x' & x' & x'  \\
						   0  & x' & x' & x'   \\
						   0  & x' & x' & x'  \end{bmatrix}$ multiply by $M_1$\\
        & $\longrightarrow$ &
		$\begin{bmatrix} x &   x        & x   & x     \\
						 0 & \boxed{x'} & x'  & x'   \\
						 0 &   0        & x'' & x''  \\
						 0 &   0        & x'' & x''  \end{bmatrix}$ multiply by $M_2$\\
        & $\longrightarrow$ &
		$\underbrace{\begin{bmatrix} x & x  &     x       & x     \\
						 0 & x' &     x'      & x'   \\
						 0 & 0  & \boxed{x''} & x''  \\
						 0 & 0  &     0       & x'''  \end{bmatrix}}_{\text{Let's call this $U$}}$ multiply by $M_3$
	\end{tabular}
\end{center}

\keyidea{
Repeatedly multiply $A$ by \textbf{elimination matrices} to make it upper triangular. 
}

\end{frame}

%%%%%%%%%%%%%%%%%

\begin{frame}{How to Construct $U$}
\begin{itemize}
    \item Repeatedly multiply $A$ from the left by elimination matrices 
        \begin{align*}
         A & \\
         M_1A  &\\
         M_2 M_1A  &\\
                  \vdots  &\\ 
         (M_{n-1} \dots M_1)A & = U
        \end{align*}
\end{itemize}

\question{
\begin{itemize}
    \item What does the product $(M_{n-1}\dots M_1)$ look like?
    \item How do you construct $L$?
\end{itemize}

}


\end{frame}


%%%%%%%%%%%%%%%%%%%%%%%%%%%%%%%%%%%%%%%%%%%%%%%%%%%%%%%%%%%%%%%%%%%%

% \begin{frame}{The Elimination Matrix}
% \define{Elimination Matrix}{ The \textbf{elimination matrix } $M_k$ is a unit lower triangular matrix which satisfies the following multiplication property
% \[
% M_k a = 
% \begin{bmatrix}
% 1 & \dots & 0 & 0 & \dots & 0\\
% \vdots & \ddots & \vdots & \vdots &  & \vdots \\
% 0 & \dots & 1 & 0 & \dots & 0\\
% 0 & \dots & -m_{k+1} & 1 & \dots & 0\\
% \vdots &  & \vdots & \vdots & \ddots & \vdots \\
% 0 & \dots & -m_{n} & 0 & \dots & 1\\
% \end{bmatrix}
% \begin{bmatrix}
% a_1 \\ \vdots \\ a_k \\ a_{k+1}\\ \vdots \\ a_n
% \end{bmatrix}
% =
% \begin{bmatrix}
% a_1 \\ \vdots \\ a_k \\ 0\\ \vdots \\ 0
% \end{bmatrix}
% \]
% }
% \begin{itemize}
% \item $M_k$ is designed to eliminate all entries of $a$ \underline{below the $k^{th}$ element}
% \item The multipliers are given by  $$m_i = \frac{a_i}{a_k}, \quad i=k+1,\dots,n$$
% \end{itemize}
% \end{frame}

\sectionframe{Some Properties}

%%%%%%%%%%%%%%%%%%%%%%%%%%%%
\begin{frame}{Formula for Elimination Matrices}
\begin{itemize}
    \item This is the $M_1$ elimination matrix because it eliminates elements under the first pivot 
        $$
        M_1=\begin{bmatrix}
    1 & 0 & 0 & 0\\
    -2 & 1 & 0 & 0\\
    -\frac{1}{2} & 0 & 1 & 0\\
    1 & 0 & 0 & 1\\
    \end{bmatrix}
        $$
    \item Notice that this matrix can be written as
            $$
        M_1=\begin{bmatrix}
    1 & 0 & 0 & 0\\
    -2 & 1 & 0 & 0\\
    -\frac{1}{2} & 0 & 1 & 0\\
    1 & 0 & 0 & 1\\
    \end{bmatrix} = \begin{bmatrix}
    1 & 0 & 0 & 0\\
    0 & 1 & 0 & 0\\
    0 & 0 & 1 & 0\\
    0 & 0 & 0 & 1\\
    \end{bmatrix}-\underbrace{\begin{bmatrix}
    0 & 0 & 0 & 0\\
    2 & 0 & 0 & 0\\
    \frac{1}{2} & 0 & 0 & 0\\
    -1 & 0 & 0 & 0\\
    \end{bmatrix}}_{m^{(1)}e_1^T},
        $$
    where
    \[
    m^{(1)}=
    \begin{bmatrix}
        0\\2\\\frac12\\-1
    \end{bmatrix},
    \qquad
    e_1^T=
    \begin{bmatrix}
        1&0&0&0
    \end{bmatrix}.
    \]
   
\end{itemize}
\end{frame}
%%%%%%%%%%%%%%%%%%%%%%%%%%%%%%%%%%%%%%%%%%%%%%%%%%%%%%%%%%%%%%%%%%%%%%%

\begin{frame}{Elimination Matrix Properties}
    \begin{block}{Elimination Matrix}
    Elimination matrices have the form
        $$ M_k = I - m^{(k)} e_k^T,$$
        where 
    \begin{itemize}
    \item $e_k$ is the $k^{th}$ column of the identity matrix $I$
    \item $m^{(k)} = [0,\dots,0,m_{k+1},\dots,m_n]'$ vector with multipliers
    \end{itemize}
    \end{block}

  \begin{block}{Property 1: Inverse}
  The inverse is $M^{-1}_k = I + m^{(k)} e_k^T$
  \end{block}
  
  \begin{block}{Property 2: Multiplication}
  Product of two elimination matrices $M_k$ and $M_j$ is a combined elimination matrix      
  \end{block}
\end{frame}

%%%%%%%%%%%%%%%%%%%%%%%%%%%%%%%%%%%%%%%%%%%%%%%%%%%%%%%%%%%%%%%%%%%%%%%%

% \begin{frame}{General Form}
% \begin{itemize}
%     \item Let $a=\begin{bmatrix} 1\\  2 \\ 4 \\ -2\end{bmatrix}$ and 
%     $I =     \begin{bmatrix}
%     1 & 0 & 0 &0\\
%     0 & 1 & 0 & 0\\
%     0 & 0 & 1 & 0\\
%     0 & 0 & 0 & 1\\
%     \end{bmatrix}$. 
%     \item For $k=2
%     $, we have  $m^{(k)}=\begin{bmatrix} 0\\  0 \\ 2 \\ -1\end{bmatrix}$ and $e_k =\begin{bmatrix} 0\\  1 \\ 0 \\ 0\end{bmatrix}$.
%     \item Thus the elimination matrix is $M_k = I - m^{(k)} e_k^T$ 
%     \begin{align*}
%     M_k &= \begin{bmatrix}
%     1 & 0 & 0 &0\\
%     0 & 1 & 0 & 0\\
%     0 & 0 & 1 & 0\\
%     0 & 0 & 0 & 1\\
%     \end{bmatrix} - \begin{bmatrix} 0\\  0 \\ 2 \\ -1\end{bmatrix} 
%     \begin{bmatrix} 0 & 1 & 0 & 0\end{bmatrix}\\
%     &=  \begin{bmatrix}
%     1 & 0 & 0 &0\\
%     0 & 1 & 0 & 0\\
%     0 & 0 & 1 & 0\\
%     0 & 0 & 0 & 1\\
%     \end{bmatrix}
%     - \begin{bmatrix} 
%     0 & 0 & 0 & 0 \\ 
%     0 & 0 & 0 & 0 \\ 
%     0 & 2 & 0 & 0 \\ 
%     0 & -1 & 0 & 0
%     \end{bmatrix}
%     = \begin{bmatrix} 
%     1 & 0 & 0 & 0 \\ 
%     0 & 1 & 0 & 0 \\ 
%     0 & -2 & 1 & 0 \\ 
%     0 & 1 & 0 & 1
%     \end{bmatrix}.
%     \end{align*}
    
% \end{itemize}
% \end{frame}

%%%%%%%%%%%%%%%%%%%%%%%%%%%%%%%%%%%%%%%%%%%%%%%%%%%%%%%%%%%%%%%%%%%%%%%%
\begin{frame}[shrink]{Inverse Property}
\begin{theorem}
If $M_k = I - m^{(k)} e_k'$ then $M_k^{-1} = I+ m^{(k)} e_k'$. 
\end{theorem}
Proof: 
\begin{align*}
    M_k M_k^{-1} &= (I - m^{(k)} e_k')(I + m^{(k)} e_k') \\
    &= II + Im^{(k)} e_k' - m^{(k)} e_k'I - m^{(k)} e_k' m^{(k)} e_k' \\
    &= I + \underbrace{m^{(k)} e_k' - m^{(k)} e_k'}_{\text{cancel}} - m^{(k)} (e_k' m^{(k)}) e_k' \\
    &= I -  m^{(k)} \underbrace{(e_k' m^{(k)})}_{\text{i.p.}=0} e_k' \\
    &= I 
\end{align*}

\keyidea{
\begin{itemize}
    \item To find the inverse simply \underline{flip the signs of the multipliers}.
    \item For instance:
        \[
        \textrm{If }\; M_1  =  
        \begin{bmatrix}
          1 & 0 & 0\\
        -2 & 1 & 0\\
        1 & 0 & 1\\
        \end{bmatrix},
         \quad \textrm{then} \;\;
         M_1^{-1}=
        \begin{bmatrix}
          1 & 0 & 0\\
        2 & 1 & 0\\
        -1 &  0 & 1\\
        \end{bmatrix}
        \]
\end{itemize}
}
\end{frame}

%%%%%%%%%%%%%%%%%%%%%%%%%%%%%%%%%%%%%%%%%%%%%%%%%%%%%%%%%%%%%%%%%%%%%%%%

\begin{frame}[shrink]{Multiplication Property}

\begin{theorem}
    Suppose $M_j$ and $M_k$ are elementary elimination matrices with $j>k$,
    then their product is $I - m^{(k)}e_k^T - m^{(j)} e_j^T$. 
\end{theorem}

    Proof: 
    \begin{align*}
      M_k M_j & = I - m^{(k)} e_k^T - m^{(j)} e_j^T + m^{(k)} e_k^T m^{(j)} e_j^T\\
    & = I - m^{(k)} e_k^T - m^{(j)} e_j^T + m^{(k)} (e_k^T m^{(j)}) e_j^T\\
    & = I - m^{(k)} e_k^T - m^{(j)} e_j^T
    \end{align*}
    
\keyidea{
\begin{itemize}
    \item To multiply $M_k M_j$ simply \underline{combine their columns}.
    \item For instance: 
        \[
        M_1 M_2 =  
        \begin{bmatrix}
          1 & 0 & 0\\
        -2 & 1 & 0\\
        1 & 0 & 1\\
        \end{bmatrix}
        \begin{bmatrix}
          1 & 0 & 0\\
        0 & 1 & 0\\
        0 & 1/2 & 1\\
        \end{bmatrix} = 
        \begin{bmatrix}
          1 & 0 & 0\\
        -2 & 1 & 0\\
        1 & 1/2 & 1\\
        \end{bmatrix}
        \]
\end{itemize}
}
\end{frame}


%%%%%%%%%%%%%%%%%%%%%%%%%%%%%%%%%%%%%%%%%%%%%%%%%%%%%%%%%%%%%%%%%%%%%%%%
\sectionframe{Constructing $L$ and $U$}
%%%%%%%%%%%%%%%%%%%%%%%%%%%%%%%%%%%%%%%%%%%%%%%%%%%%%%%%%%%%%%%%%%%%%%%%

\begin{frame}{Objective}
The goal is write $A$ as $A = LU$, where 
\begin{equation*}
    L = \begin{bmatrix}   1       &     0      & \cdots &             & 0 \\
                       \ell_{2,1} &     1      &   0    &             & 0 \\
                       \ell_{3,1} & \ell_{3,2} &   1    &             & 0 \\
                          \vdots  &  \vdots    &        & \ddots      & \vdots  \\
                       \ell_{n,1} & \ell_{n,2} & \cdots & \ell_{n,n-1} & 1
        \end{bmatrix} \quad 
    U = \begin{bmatrix} u_{1,1} & u_{1,2} & u_{1,3} & \cdots & u_{1,n} \\
                          0     & u_{2,2} & u_{2,3} & \cdots & u_{2,n} \\
                          0     &   0     & \ddots  & \ddots & \vdots \\
                         \vdots & \vdots  &         &        & u_{n-1,n} \\
                          0     &   0     &         &        & u_{n,n}
        \end{bmatrix}
\end{equation*}

\onslide<2>{
\question{Don't we already know a way to make $A$ upper triangular?}
}

\end{frame}



%%%%%%%%%%%%%%%%%%%%%%%%%%%%%%%%%%%%%%%%%%%%%%%%%%%%%%%%%%%%%%%%%%%%%%%%

\begin{frame}{LU Factorization}{Part 1}
\begin{itemize}
    \item[1.] Repeatedly multiply $A$ from the left by elimination matrices 
        \begin{align*}
         A & \\
         M_1A  &\\
         M_2 M_1A  &\\
                  \vdots  &\\ 
         \underbrace{(M_{n-1} \dots M_1)A}_{\textrm{upper triangular}} & = U
        \end{align*}
    \item[2.] In the end, the product will be upper triangular
\end{itemize}



\end{frame}

%%%%%%%%%%%%%%%%%%%%%%%%%%%%%%%%%%%%%%%%%%%%%%%%%%%%%%%%%%%%%%%%%%%%%%%%%

\begin{frame}{LU Factorization}{Part 2}

\begin{itemize}
    \item[3.] Multiply both sides from the left by
    $(M_{n-1}\cdots M_1)^{-1}$:
    \begin{align*}
        (M_{n-1}\cdots M_1)A &= U,\\
        A &= (M_{n-1}\cdots M_1)^{-1}U.
    \end{align*}

    \item[4.] Reverse the order when taking the inverse:
    \[
        A=M_1^{-1}\cdots M_{n-1}^{-1}U.
    \]
\end{itemize}

\begin{block}{Inverse of a Product}
For square nonsingular matrices $P$ and $Q$,
\[
    (PQ)^{-1}=Q^{-1}P^{-1}.
\]
\end{block}

\end{frame}

%%%%%%%%%%%%%%%%%%%%%%%%%%%%%%%%%%%%%%%%%%%%%%%%%%%%%%%%%%%%%%%%%%%%%%%%%

\begin{frame}{LU Factorization}{Part 3}

\begin{itemize}
        
    \item[5.] Use the two \textbf{multiplication and inverse properties}: 

    $$(M_1^{-1} \dots M_{n-1}^{-1}) = I + \sum_{k=1}^{n-1} m^{(k)}e_k^T = L$$
    Invert each elimination matrix by changing the signs
of its off-diagonal entries, then combine the columns.

    \item[6.] We are done with the LU factorization 
           \begin{align*}
            A & = (M_1^{-1} \dots M_{n-1}^{-1})U\\
            A & = LU
        \end{align*}
        
\keyidea{
To construct $L$:
\begin{itemize}
    \item Put ones on the diagonal.
    \item Store each elimination multiplier
    $m_{ik}=a_{ik}/a_{kk}$ in position $(i,k)$.
    \item These entries have the opposite signs
    to the corresponding entries in $M_k$.
\end{itemize}
}

\end{itemize}
\end{frame}

%%%%%%%%%%%%%%%%%%%%%%%%%%%%%%%%%%%%%%%%%%%%%%%%%%%%%%%%%%%%%%%%%%%%%%%%%
\begin{frame}{Factorization Example}{Part 1}

\begin{example}
    Find the $LU$ factorization of $A  =
\begin{bmatrix}
1 & -2 & 3\\
-1 & 4 & -2\\
2 & 2 & 13\\
\end{bmatrix}$.
\end{example}

\begin{itemize}
    \item[1.] Construct elimination matrix $M_1$. Left-multiply it to $A$.
\begin{align*}
M_1 A = 
\begin{bmatrix}
1 & 0 & 0\\
1 & 1 & 0\\
-2 & 0 & 1\\
\end{bmatrix}
\begin{bmatrix}
\boxed{1} & -2 & 3\\
-1 & 4 & -2\\
2 & 2 & 13\\
\end{bmatrix} = \begin{bmatrix}
1 & -2 & 3\\
0 & 2 & 1\\
0 & 6 & 7\\
\end{bmatrix}
\end{align*}

\item[2.] Construct elimination matrix $M_2$. Left-multiply it. \underline{$U$ is complete}.
\begin{align*}
M_2 M_1 A = 
\begin{bmatrix}
1 & 0 & 0\\
0 & 1 & 0\\
0 & -3 & 1\\
\end{bmatrix}
\begin{bmatrix}
1 & -2 & 3\\
0 & \boxed{2} & 1\\
0 & 6 & 7\\
\end{bmatrix} = \begin{bmatrix}
1 & -2 & 3\\
0 & 2 & 1\\
0 & 0 & 4\\
\end{bmatrix} =: U
\end{align*}
\end{itemize}
\end{frame}
%%%%%%%%%%%%%%%%%%%%%%%%%%%%%%%%%%%%%%%%%%%%%%%%%%%%%%%%%%%%%%%%%%%%%%%%%
\begin{frame}{Factorization Example}{Part 2}

\begin{itemize}

\item[3.] Derive a formula for $L$ based on the equation $A=LU$
\begin{align*}
    M_2 M_1 A &= U\\
    A &= M_1^{-1} M_2^{-1} U\\ 
      \therefore L &= M_1^{-1} M_2^{-1}
\end{align*}

\item[4.] Construct $L$ using the properties of elimination matrices:
$$
M_1 = \begin{bmatrix}
1 & 0 & 0\\
1 & 1 & 0\\
-2 & 0 & 1\\
\end{bmatrix} \implies 
M_1^{-1} = \begin{bmatrix}
1 & 0 & 0\\
-1 & 1 & 0\\
2 & 0 & 1\\
\end{bmatrix}\quad\text{just flip sign}
$$
$$
M_2 = \begin{bmatrix}
1 & 0 & 0\\
0 & 1 & 0\\
0 & -3 & 1\\
\end{bmatrix} \implies 
M_2^{-1} = \begin{bmatrix}
1 & 0 & 0\\
0 & 1 & 0\\
0 & 3 & 1\\
\end{bmatrix} \quad\text{just flip sign}
$$
$$
L = M_1^{-1}M_2^{-1}= \begin{bmatrix}
1 & 0 & 0\\
-1 & 1 & 0\\
2 & 3 & 1\\
\end{bmatrix} \quad\text{just combine columns}
$$
\end{itemize}
\end{frame}
%%%%%%%%%%%%%%%%%%%%%%%%%%%%%%%%%%%%%%%%%%%%%%%%%%%%%%%%%%%%%%%%%%%%
\begin{frame}{Summary}
\begin{itemize}
    \item The $LU$ factorization of the matrix is 
    $$
    \begin{bmatrix}
1 & -2 & 3\\
-1 & 4 & -2\\
2 & 2 & 13\\
\end{bmatrix} = 
\begin{bmatrix}
1 & 0 & 0\\
-1 & 1 & 0\\
2 & 3 & 1\\
\end{bmatrix}
\begin{bmatrix}
1 & -2 & 3\\
0 & 2 & 1\\
0 & 0 & 4\\
\end{bmatrix}
    $$
    \item You may multiply $LU$ to verify that it equals $A$\\[3mm]

    \item Main steps:
    \begin{enumerate}
        \item Perform GE with elimination matrices
        \item Repeat until the result is upper triangular - that is $U$
        \item Construct $L$ by combining all eliminations but with sign-flipped multipliers
    \end{enumerate}
\end{itemize}

    
\end{frame}

%%%%%%%%%%%%%%%%%%%%%%%%%%%%%%%%%%%%%%%%%%%%%%%%%%%%%%%%%%%%%%%%%%%%
\begin{frame}{Coded Example}

\begin{block}{Matlab: \texttt{matrix\_form\_LU.m}}
    
The code demonstrates the $LU$ factorization of $A  =
\begin{bmatrix}
1 & 2 & 3\\
2 & 5 & 1\\
3 & 4 & 1\\
\end{bmatrix}$.
\end{block}

\begin{itemize}
    \item Run one section at a time
    \item Understand how row operations are coded into elimination matrices
\end{itemize}


$$ 
L =
\begin{bmatrix}
1 & 0 & 0\\
2 & 1 & 0\\
3 & -2 & 1\\
\end{bmatrix} \quad 
U = 
\begin{bmatrix}
1 & 2 & 3\\
0 & 1 & -5\\
0 & 0 & -18\\
\end{bmatrix}
$$
\end{frame}

%%%%%%%%%%%%%%%%%%%%%%%%%%%%%%%%%%%%%%%%%%%%%%%%%%%%%%%%%%%%%%%%%%%%%%%%%

\begin{frame}[fragile]{LU versus GE}
\begin{columns}
    \begin{column}{0.45\textwidth}
\begin{lstlisting}[mathescape]
given $A$
                           
for $k=1\dots n-1$          
    for $i = k+1\ldots n$      
        $xmult = a_{ik}/a_{kk}$
        $a_{ik} = xmult$      
        for $j = k+1\ldots n$      
            $a_{ij} = a_{ij} - (xmult) a_{kj}$
        end                      
        % doesn't affect b
    end                        
end                          
\end{lstlisting}
\begin{center}
    LU factorization 
\end{center}
    \end{column}
    \begin{column}{0.45\textwidth}
\begin{lstlisting}[mathescape,label=algo:fwdelim]
given $A$, $b$
                               
for $k=1\dots n-1$          
    for $i = k+1\ldots n$      
        $xmult = a_{ik}/a_{kk}$
        $a_{ik} = 0$     
        for $j = k+1\ldots n$      
            $a_{ij} = a_{ij} - (xmult) a_{kj}$
        end                      
        $b_{i} = b_{i} - (xmult) b_{k}$
    end                        
end                          
\end{lstlisting}
\begin{center}
    GE
\end{center}
    \end{column}
\end{columns}
\vspace{0.5cm}
\begin{center}
They are basically the same algorithm.
\end{center}
\end{frame}

%%%%%%%%%%%%%%%%%%%%%%%%%%%
\begin{frame}{LU Factorization}
    \method{LU Factorization}{Square matrix $A$}{Unit lower $L$ and upper triangular $U$, such that $LU=A$}
    {Factorization of $A$ comes naturally from GE.}

\vspace{3mm}
Recap:
\begin{itemize}
    \item $U$ is naturally constructed as part of Gaussian elimination.
        $$U = (M_{n-1}\cdots M_1)A$$
    \item $L$ is a byproduct of GE adding no extra cost (just save the multipliers) 
        $$ L = M_1^{-1} \cdots M_{n-1}^{-1}$$
\end{itemize}
\keyidea{
LU factorization performs elimination and saves the
multipliers: $\mO(n^3)$ FLOPs.

Once the factors are available, an LU solve uses
forward and backward substitution: $\mO(n^2)$ FLOPs
per right-hand side.
}
\end{frame}
%%%%%%%%%%%%%%%%%%%%%%%%%%%%%%%
\begin{frame}{Questions}

\begin{enumerate}
    \item Do all matrices have an $LU$ factorization? 
\vfill
    \item Does $LU$ factorization need pivoting like GE?
\vfill
    \item Are there other useful factorizations? 
\vfill
\end{enumerate}
    
\end{frame}


%%%%%%%%%%%%%%%%%%%%%%%%%%%%%%%

\end{document}